\documentclass[a4paper, 11pt]{article}

%%% SST LAB PROTOCOLL PREAMBLE
%%% 2019
%%%%%%%%%%%%%%%%%%%%%%%%%%%%%%%


%%% PACKAGES
%%%%%%%%%%%%%%%%%%%%%%%%%%%

\usepackage[ngerman]{babel}

\usepackage[utf8]{inputenc}
\usepackage{amsmath}
\usepackage{pgfplots}
\usepackage{tikz}
\usepackage[many]{tcolorbox}
\usepackage{graphicx}
\graphicspath{ {./graphics/} }
\usepackage{pdfpages}
\usepackage{dashrule}
\usepackage{float}
\usepackage{siunitx}
\usepackage{trfsigns}
\usepackage{booktabs}
\usepackage[european]{circuitikz}
\usepackage{tcolorbox}

%%% DOCUMENT GEOMETRY
%%%%%%%%%%%%%%%%%%%%%%%%%%%

\usepackage{geometry}
\geometry{
 a4paper,
 total={0.6180339887498948\paperwidth,0.6180339887498948\paperheight},
 top = 0.1458980337503154\paperheight,
 bottom = 0.1458980337503154\paperheight
 }
\setlength{\jot}{0.013155617496424828\paperheight}
\linespread{1.1458980337503154}

\setlength{\parskip}{0.013155617496424828\paperheight} % paragraph spacing


%%% COLORS
%%%%%%%%%%%%%%%%%%%%%%%%%%%

\definecolor{red1}{HTML}{f38181}
\definecolor{yellow1}{HTML}{fce38a}
\definecolor{green1}{HTML}{95e1d3}
\definecolor{blue1}{HTML}{66bfbf}
\definecolor{hsblue}{HTML}{00b1db}
\definecolor{hsgrey}{HTML}{afafaf}

%%% CONSTANTS
%%%%%%%%%%%%%%%%%%%%%%%%%%%
\newlength{\smallvert}
\setlength{\smallvert}{0.0131556\paperheight}


%%% COMMANDS
%%%%%%%%%%%%%%%%%%%%%%%%%%%

% differential d
\newcommand*\dif{\mathop{}\!\mathrm{d}}

% horizontal line
\newcommand{\holine}[1]{
  	\begin{center}
	  	\noindent{\color{hsgrey}\hdashrule[0ex]{#1}{1pt}{3mm}}\\%[0.0131556\paperheight]
  	\end{center}
}

% mini section
\newcommand{\minisec}[1]{ \noindent\underline{\textit {#1} } \\}

% quick function plot
\newcommand{\plotfun}[3]{
  \vspace{0.021286\paperheight}
  \begin{center}
    \begin{tikzpicture}
      \begin{axis}[
        axis x line=center,
        axis y line=center,
        ]
        \addplot[draw=red1][domain=#2:#3]{#1};
      \end{axis}
    \end{tikzpicture}
  \end{center}
}

% box for notes
\newcommand{\notebox}[1]{

\tcbset{colback=white,colframe=green1!100!black,title=Note!,width=0.618\paperwidth,arc=0pt}

 \begin{center}
  \begin{tcolorbox}[]
   #1 
  \end{tcolorbox}
 
 \end{center} 
 
}

% box for equation
\newcommand{\eqbox}[2]{
	
	\tcbset{colback=white,colframe=green1!100!black,title=,width=#2,arc=0pt}
	
	\begin{center}
		\begin{tcolorbox}[ams align*]
				#1
		\end{tcolorbox}
		
	\end{center} 
	
}
% END OF PREAMBLE

%%%%%%%%%%%%%%%%%%%%%%%%%%%%%%%%%%%%%

\begin{document}

%%%%%%%%%%%%%%%%%%%%%%%%%%%%%%%%%%%%%
  \includepdf{./titlepage/titlepage.pdf}
  \clearpage
  \setcounter{page}{1}
%%%%%%%%%%%%%%%%%%%%%%%%%%%%%%%%%%%%%

  \section{RS232-Datenübertragung}
  \label{sec:1}
  Die Datenübetragung mittels RS232 (Recommended Standard 232) erfolgt asynchron
  und seriell. Das Übertragungsformat und die Übertragungsgeschwindigkeit werden
  nicht im RS232 definiert sondern durch die Interface-Hardware
  bestimmt, in der Regel UART. Der Standard legt die Spannungspegel zur
  Übertragung fest. Diese liegen im Bereich $-3.3 \, \si{\volt} ... -15 \,
  \si{\volt}$ zum Senden einer 1 und $+3.3 \si{\volt} ... +15 \, \si{\volt}$ zum
  Senden einer 0.

  \section{Datensynchronisation zwischen Sender und Empfänger}
  Zur Synchronisation zwischen Sender und Empfänger dienen die Start- und
  Stopp-Bits am Anfang bzw. Ende jedes Übertragungsframes. Da die
  Übertragungsrate (Baudrate) Sender und Empfänger bekannt ist, wird dadurch
  kein separates Taktsignal benötigt.

  \section{Einstellung der Parameter, Schnittstelle USCI\_A1, }
  
  \begin{itemize}
  \item[]{
      \inlinecode{UCA1CTL0}
      \begin{itemize}
        \item []\inlinecode{UCMODE1} \commentGray{UART Mode}
        \item []\inlinecode{UCSPB} \commentGray{1 Stop-Bit}
        \item []\inlinecode{UCSYNC} \commentGray{0...Asynchroner Modus}
      \end{itemize}
    }

   \item[]{
      \inlinecode{UCA1CTL1}
      \begin{itemize}
        \item []\inlinecode{UCSSEL1} \commentGray{BRCLK Taktquelle}
        \item []\inlinecode{UCSWRST} 
      \end{itemize}
    }

   \item[]{
      \inlinecode{UCA1BR0} \commentGray{Taktteiler $N$ für Baudrate $N = \frac{f_{BRCLK}}{\mathrm{Baudrate}}$}
    }

   \item[]{
      \inlinecode{UCA1RXBUF}
    }

   \item[]{
      \inlinecode{UC1IE}
      \inlinecode{UC1IFG} \commentGray{Interrupts}
    }


  \end{itemize}

  \section{Erhöhung der Übertragungssicherheit}
  Die UART-Schnittstelle kann konfiguriert werden, um am Ende jedes Frames ein
  Paritätsbit zu senden/empfangen. Die Notation für das Übertragungsformat ist
  dann z.B. 8/{\color{cyan5}{E}}/1 oder 8/{\color{orange5}{O}}/1 für gerade ({\color{cyan5}{E}}ven)
  bzw. ungerade ({\color{orange5}{O}}dd) Parität.

  \section{MAX3227E}
  Zur Konvertierung der TTL/CMOS-Pegel ($0 \, \si{\volt}...\pm 3.3 \,
  \si{\volt}/ \pm 5\,\si{\volt}$) des Mikrocontrollers auf die RS232-Pegel
  (siehe \ref{sec:1}) wird ein Pegelwandler benötigt, in diesem Fall der MAX3227E.
  
\end{document}
% Local Variables:
% TeX-engine: luatex
% End:
